\begin{titlepage}
\color{bleu}\sffamily
{\large MATHS EXPLORER\par}
\vspace{22mm}
{\fontsize{35}{40}\selectfont\bfseries Mathématiques\par}
\vspace{6mm}
{\LARGE Première\par spécialité\par}
\vspace{12mm}
{\Large 2026--2027\par}
\vspace{10mm}
\begin{tikzpicture}[scale=.8]\draw[vert,very thick,domain=-1.7:1.7,samples=80]plot(\x,{\x*\x});\draw[bleu,thick](-2,0)--(2,0);\draw[bleu,thick](0,-.3)--(0,3.2);\draw[ocre,thick](-.5,-2)--(2,3);\end{tikzpicture}
\vfill
{\Large\bfseries\ifprof Édition professeur\else Édition élève\fi\par}
\vspace{4mm}
Cours développés, parcours progressifs et préparation à l’épreuve anticipée.\par
Dix chapitres dans l’ordre de la progression de l’établissement.\par
\vspace{8mm}
{\small Version éditoriale révisée du 10 septembre 2026.\par À utiliser avec CQFD Mathématiques 1re, Bordas, 2019.}
\end{titlepage}
\chapter*{Utiliser ce livret}
\addcontentsline{toc}{chapter}{Utiliser ce livret}
Le cours contient des exemples entièrement résolus et les démonstrations à connaître. Les diagnostics, applications, exercices, QCM et entraînements sont à traiter sur un cahier. Le professeur dispose des corrections et de variantes réservées. Le livret suit la progression fournie ; les numéros du manuel sont conservés dans les références et ne se confondent pas avec les numéros propres à ce livret.

Chaque chapitre commence par six questions de prérequis et une activité de recherche. Le parcours conduit ensuite des exercices guidés au choix d’une méthode, puis au transfert. L’exercice 8 est prévu pour une reprise différée. Les exercices de réserve ne constituent pas une liste de devoirs à terminer intégralement.

\begin{repere}{Comment travailler}
Chercher d’abord sans solution. Écrire les données, le domaine et la question. Après une aide ou une correction, refermer le cours et reconstruire une étape. La semaine suivante, refaire une tâche courte sans notes. Une réponse fausse accompagnée d’un essai exploitable permet au professeur de cibler l’aide.
\end{repere}

\textbf{Rédiger.} Une égalité est justifiée par une définition, une propriété ou un calcul. Une conclusion répond à la question, avec une unité si nécessaire. « Conjecturer » signifie proposer une réponse plausible ; « démontrer » signifie établir sa validité sous des hypothèses explicites.

\textbf{Outils.} Les exercices numériques peuvent utiliser Python quand cela est demandé. Pour les automatismes et les sujets de préparation à l’EAM, la calculatrice est interdite. Pour un travail personnel, conserver le premier essai et signaler les aides reçues ; pouvoir expliquer sa solution sans outil est un objectif central.

\ifprof\paragraph{Retour des évaluations.} Conserver trois indicateurs : exactitude, choix de méthode, justification. Un QCM réussi ne certifie pas la rédaction. Refaire une tâche incorrecte sans cours, puis à une semaine et au créneau de consolidation.\fi

\textbf{Notations.} Les intervalles utilisent les conventions françaises. Une probabilité conditionnelle peut être notée $P_A(B)$ ou $P(B\mid A)$. Les figures géométriques sont en repère orthonormé lorsque les calculs l’exigent. Les figures de cours illustrent une propriété ; elles ne remplacent pas sa preuve.

\chapter*{Cadre de l’année et de l’EAM}
\addcontentsline{toc}{chapter}{Cadre de l’année et de l’EAM}
Le programme retenu est celui de l’arrêté du 26 février 2026 (NOR MENE2602917A), publié au BO no 14 du 2 avril 2026. Son article 2 le rend applicable à la rentrée 2026--2027. Il remplace l’annexe du programme de première de 2019. Le nouveau programme de terminale prévu pour la rentrée 2027--2028 n’est pas celui de ce livret.

Les élèves de première 2026--2027 passent l’épreuve anticipée en 2027, au titre du baccalauréat 2028. La note de service du 10 juin 2025 (NOR MENE2515469N) définit une épreuve de deux heures, coefficient 2, sans calculatrice : QCM d’automatismes sur 6 points, puis deux ou trois exercices indépendants sur 14 points. La liste transitoire d’automatismes portant explicitement sur 2025--2026 ne remplace pas les automatismes du programme 2026 applicable à cette cohorte.

\textbf{Format pédagogique du livret.} Chaque chapitre contient deux QCM de cinq questions. Les sujets annuels utilisent douze questions à 0,5 point, suivies de trois exercices totalisant 14 points : c’est un choix de préparation compatible avec la structure réglementaire, pas une affirmation sur le nombre de questions du futur sujet officiel. Le sujet métropolitain de juin 2026 comporte par exemple huit questions d’automatismes.

\textbf{Calendrier.} Le groupe de référence dispose de créneaux de 2 h les mardis et jeudis, de 9 h 45 à 11 h 45. Le DS commun est annoncé le 14 novembre ; les blancs du 14 au 16 janvier et du 1er au 3 avril. Leurs horaires précis ne sont pas fournis. Les sujets proposés durent 2 h ; cette durée reste une proposition pour les évaluations communes. La date EAM de Madagascar n’est pas déduite du calendrier métropolitain.

\textbf{Références officielles.}\par
\href{https://www.education.gouv.fr/bo/2026/Hebdo14/MENE2602917A}{Programme de première, BO du 2 avril 2026}.\par
\href{https://www.education.gouv.fr/bo/2025/Hebdo24/MENE2515469N}{Définition de l’épreuve, BO du 12 juin 2025}.\par
\href{https://eduscol.education.gouv.fr/5688/epreuve-anticipee-de-mathematiques-aux-baccalaureats-general-et-technologique}{Éduscol : présentation et sujets zéro}.\par
\href{https://www.apmep.fr/Epreuve-anticipee-2026}{APMEP : répertoire des annales 2026, plusieurs centres}.
\tableofcontents
\clearpage
